\section{External Constraints}
All space agencies must adhere to international space law and regulations, using these well-defined guidelines as a foundational basis for conceptual proposals and recognizing them as essential constraints. For the subsystem design, the following constraints are considered:
\begin{itemize}
    \item Signed in 1967, by Outer Space Treaty (OST), Responsibility, and Liability: \textit{states are liable for any damages caused by their space objects} \cite{un1967ost}. 
    \item Under Article IX of the OST, Environmental Protection and Avoidance of Harmful Contamination: \textit{Satellites should be designed to minimize debris generation, preventing environmental contamination} \cite{un1972liability}.
    \item Stated by The Inter-Agency Space Debris Coordination Committee (IADC) and United Nations (UN), Space Debris Mitigation Guidelines \cite{un1967environment}:
    \begin{itemize}
        \item \textit{Satellites in low Earth orbit (LEO) should be deorbited within 25 years of mission end}.
        \item \textit{Satellite operators must monitor and avoid potential collisions with other objects.}
        \item \textit{Use of materials that minimize fragmentation risk}
    \end{itemize}
    \item Outlined by International Telecommunication Union (ITU), Radio Frequency Allocation: \textit{Satellite operators must coordinate with the ITU to assign specific radio frequencies, and systems must be designed to prevent harmful interference with other satellites} \cite{itu2008radio}.
    \item The OST encourages: \textit{to share data from scientific missions and cooperate to improve safety and sustainability in space }\cite{un1967ost}.

\end{itemize}